\section{Metodologia}\label{sec:metodologia-revisada}
O estudo articula desenvolvimento de um artefato e análise de evidências documentais e técnicas. As etapas de ouvir e interpretar o contexto, criar e prototipar, implementar e testar seguem a organização do modelo institucional fornecido e são orientadas pelo Design Thinking e pelo Double Diamond \citep{ideo-designthinking,designcouncil}. Não se afirma execução integral de oficinas formais dessas abordagens. A Tabela~\ref{tab:procedimentos-revisados} separa procedimentos realizados de avaliação futura. O histórico detalhado e os instrumentos permanecem no Apêndice~\ref{ap:percurso-metodologico}; os limites de inferência são reunidos na Seção~\ref{sec:validade}.
\begingroup\small\SingleSpacing\setlength{\parskip}{0pt}\setlength{\tabcolsep}{4pt}
\begin{longtable}{>{\raggedright\arraybackslash}p{2.6cm}p{5.7cm}p{6.4cm}}
\caption{Procedimentos, evidências e alcance das conclusões.}\label{tab:procedimentos-revisados}\\\hline
Procedimento & Método e evidência & Alcance e situação\\\hline\endfirsthead
\hline Procedimento & Método e evidência & Alcance e situação\\\hline\endhead
\hline\endfoot
Necessidades e entrevista & Síntese temática da conversa remota com Osmar Junior; transcrição identificada em 20/8/2026, 2.861 s; intervalos na Tabela~\ref{tab:entrevista-osmar}. & Realizado. Perspectiva profissional da gestão; não representa amostra de passageiros, parceria ou autorização de implantação.\\[5pt]
Pesquisa documental & Literatura sobre informação e IHC; registros públicos, e-mail de 2/9 e resposta sobre Santos, com procedência discriminada. & Realizada. Fundamenta recorte e requisitos; relatos e fontes secundárias não viram cadastros operacionais validados.\\[5pt]
Soluções existentes & Inspeção de capturas de QRCodeAdesivo, Olho Vivo e Cittamobi; análise de favoritos e descrição indireta do fluxo em Santos. & Realizada no acervo disponível. Tabela~\ref{tab:comparacao-servicos}; não é experimento comparativo com usuários.\\[5pt]
Desenvolvimento e prototipação & Mockups, Figma e aplicação Django/React; inspeção de código, requisitos retrospectivos e versões Git. & Implementação e iterações documentadas. Mockup não prova funcionamento; código não prova benefício de uso.\\[5pt]
Integração e auditoria & Importação GTFS; recorte geográfico com malha IBGE 2024; contagem de identificadores, viagens e dias de serviço; consultas temporalmente verificadas. & Realizadas; 337 paradas ARTESP e 16 linhas/variantes em Mauá. Cadastro e programação, não observação da operação em campo.\\[5pt]
Testes funcionais & Testes unitários, integração e navegador registrados por revisão; capturas automatizadas. Apêndice~\ref{ap:engenharia} e registros de 29/9 e 29/10. & Execuções históricas documentadas. Nesta revisão editorial não se repetiram testes do aplicativo; não substituem avaliação de usabilidade.\\[5pt]
Benchmarks & Repetição de 3.264 consultas por versão em 1/11, 32 processos, chamadas externas bloqueadas; CSV individual, 1.637 sucessos pareados. & Medição local até a primeira rota; gráficos da Seção~\ref{sec:desempenho-revisado}. Não mede latência de ponta a ponta ou causa isolada da variação.\\[5pt]
Avaliação profissional & Registros da Alameda Solutions de 21 e 25/10, versões e documento confirmado de 26/10; Apêndice~\ref{ap:feedback-alameda}. & Apreciação qualitativa. Não se atribuem ao avaliador testes internos nem alterações de novembro.\\[5pt]
Navegação exploratória & Relatos de usuários convidados e lista de ajustes; devolutivas de outubro e verificações técnicas da Tabela~\ref{tab:feedback-ajustes}. & Realizada com alcance exploratório. Não há amostra sistemática documentada, taxa de sucesso ou mensuração representativa de esforço.\\[5pt]
Avaliação com passageiros & Protocolo de tarefas QR versus entrada manual, compreensão, tempo, erros, ajuda e percepção de facilidade; instrumentos preservados. & Futura. Participantes, coleta e resultados ainda não apresentados. Leitura física do QR exige evidência própria, conforme protocolo do Apêndice~\ref{ap:protocolo-qr}.\\
\end{longtable}
\fonte{Elaboração do grupo com base no acervo, nas versões e nos registros de execução descritos nos apêndices (revisão em 2 nov. 2026).}
\endgroup
O problema de projeto é ajudar a pessoa no ponto a reconhecer alternativas para chegar ao destino sem conhecer previamente as linhas ou a identificação da parada. Na etapa de ouvir, a entrevista e a análise documental delimitaram esse problema. Na etapa de criar, o grupo comparou formas de acesso e organização da informação; a presença de destino e conexões no Cittamobi afastou a alegação de novidade dessas funções. Na etapa de implementar, o cenário demonstrativo de Mauá foi seguido pelo planejamento com dados SPTrans e pela ampliação ARTESP, com suas fontes separadas.

A auditoria ARTESP verificou a malha municipal e os identificadores antes de contar viagens, evitando somar uma mesma viagem a cada parada. Os instantes de 3, 7 e 8 de novembro são datas de serviço consultadas, não datas de observação da circulação. O benchmark de 1 de novembro antecede essa ampliação. As comparações usam os mesmos pares que receberam rota nas duas versões; os demais estados são contabilizados separadamente. Os procedimentos de reprodução e os arquivos utilizados estão nos Apêndices~\ref{ap:importacao} e~\ref{ap:rotas}.

A interpretação relaciona evidência a requisito e limita a conclusão ao procedimento executado. A apreciação profissional e os relatos exploratórios orientaram correções, enquanto os testes internos conferiram comportamentos. A eficácia, a eficiência e a satisfação dos passageiros permanecem objetos da avaliação futura \citep{iso924111}; não são deduzidas de implementação, quantidade de testes ou número de cliques.
